\documentclass[11pt,a4paper]{article}
\usepackage{curso}
\begin{document}
\portada{Lección}{11.5}{Agentes de investigación con modelos de lenguaje}

\noindent\textbf{Unidad 11.}\quad\textbf{Objetivos:} GRPO; recompensas de ejecución y rúbrica con un modelo juez; cómo se aprenden las trampas; lectura crítica de
Thermo-RSI.
\textbf{Prerrequisitos:} 11.4.

\section{GRPO}
Shao et al.\ (2024, DeepSeekMath): para cada tarea se genera un grupo de $G$ intentos y $A_i=(r_i-\mathrm{media}(r))/\mathrm{desv}(r)$. Se refuerzan los intentos mejores
que la media del grupo; no hace falta red de valor. Ejemplo: $r=(1;\,0{,}4;\,0;\,0{,}7)\Rightarrow A=(+1{,}28;\,-0{,}34;\,-1{,}42;\,+0{,}47)$.

\section{Ejecución y juez \evP}
Recompensa de Thermo-RSI: $r=0{,}7\,\text{exec}+0{,}3\,\text{rúbrica}$; exec determinista ($0{,}15$ corre $+\ 0{,}85$ reproduce la métrica de referencia); rúbrica
calificada por Nemotron 3 Super 120B con criterios escritos de antemano. La ejecución pesa más por ser ``la señal más difícil de engañar''.

\section{Cómo se aprenden las trampas \evS}
GRPO de juguete sobre las cinco propuestas de 11.4 (política softmax, $G=8$, 300 pasos, 200 semillas; \texttt{datos/grpo\_juguete.py}). Probabilidad final:
\begin{center}\small\begin{tabular}{lccccc}\toprule
recompensa & paralelo (mal) & 3 barridos & correcta & error & trampa\\\midrule
completa & 0{,}001 & 0{,}000 & 0{,}998 & 0{,}000 & 0{,}001\\
solo ejecución & 0{,}328 & 0{,}001 & 0{,}318 & 0{,}001 & 0{,}351\\\bottomrule
\end{tabular}\end{center}
Sin rúbrica, dos tercios de la probabilidad van a un método incorrecto o a la trampa: el agente aprende lo que la recompensa premia.

\section{Resultados de Thermo-RSI \evP}
Neagoe, Verdon y Morton (Extropic con Prime Intellect, 1 de octubre de 2026): Qwen3.6-35B-A3B (3B activos), 100 pasos de GRPO ($\sim25$ h), $\sim50$ tareas que reproducen
experimentos clásicos (máquinas de Boltzmann, vigilia--sueño). Recompensa en tareas no vistas: $0{,}127\to0{,}361$ ($\approx2{,}8\times$). Mismo conjunto, sin ajuste:
GPT-5.5 $0{,}747$, Claude Opus 4.8 $0{,}464$, Qwen3.5 ajustado $0{,}356$, Qwen3-30B-Instruct $0{,}109$, GLM-5.2 $0{,}073$.

\textbf{Lectura crítica.} Demostrado: mejora al reproducir experimentos conocidos. Objetivo declarado, no demostrado: descubrir algoritmos nuevos y coevolucionar con el
hardware; las tareas aún no corren sobre el simulador de TSU. No se reportan observaciones de \emph{reward hacking}.

\begin{aplicacion}
GRPO sobre tareas verificables; recompensa dominada por la ejecución; juez con rúbrica fija; auditoría de ejecuciones. Laboratorio L30.
\end{aplicacion}

\begin{error}
\begin{itemize}[nosep]
\item Premiar solo la métrica. \quad Juez sin criterios fijados de antemano.
\item Dar por demostrada la automejora recursiva. \quad Comparar puntuaciones de conjuntos distintos.
\end{itemize}
\end{error}

\begin{thebibliography}{9}\small
\bibitem{rsi} A. Neagoe, G. Verdon y S. Morton (Extropic, con Prime Intellect), First sparks of thermodynamic recursive intelligence, 1 de octubre de 2026.
\bibitem{grpo} Z. Shao et al., DeepSeekMath, arXiv:2402.03300 (2024).
\bibitem{judge} L. Zheng et al., Judging LLM-as-a-judge with MT-Bench and Chatbot Arena, \emph{NeurIPS} (2023).
\end{thebibliography}
\end{document}
